\section{Discussion}

\parab{Additional parallelism.} While \sys primarily focuses on Expert Parallelism and Sequence Parallelism, we acknowledge other strategies not explicitly covered in our evaluation. Tensor Parallelism, widely adopted in serving, is typically confined to intra-node fabrics (e.g., NVLink) without directly contending for inter-node bandwidth. Similarly, we do not explicitly discuss Pipeline Parallelism, as the prefill stage is usually latency-sensitive (requiring low or zero PP degrees~\cite{distserve}), and its communication volume is orders of magnitude smaller than KV cache movement or collective communications~\cite{liao2025mixnet}. Should these flows appear on the shared network, they would be identified as Stage~2 traffic and remain compatible with our analysis. This observation extends to emerging parallelism strategies, which may introduce novel communication patterns with enhanced overlapping capabilities, yet still block dependent computation if delayed.

\parab{Hybrid model deployment.}
In practice, service providers often deploy a tiered model portfolio, consisting of a single large flagship model alongside multiple smaller models serving diverse workloads~\cite{openai2024gpt4}.
In such settings, network contention primarily arises from concurrent KV cache transfers (Stage~1 and Stage~3) across different models, as collective communication is typically confined within isolated racks or dedicated pods~\cite{deepseektech}.
While \sys{} is not explicitly evaluated under multi-model deployment, its scheduling principle remains applicable to KV cache traffic on shared networks.
While promising, the holistic orchestration of such multi-model contention landscapes remains an under-explored design space.

\parab{Applicability to scale-up fabrics.} While \sys focuses on scale-out networks, its design principles could be extended to scale-up domains (e.g., NVL~\cite{nvidia_gb200_nvl72}, UB~\cite{huawei384},etc.). Contention between collective and KV Cache movement remains when multiple serving instance deployed together. \sys's scheduling policy can be effectively adapted to such architectures by leveraging hardware-level isolation mechanisms (e.g., virtual channels or traffic classes), which could be exploited to enable differentiated prioritization.

\parab{Compatibility with GPU-initiated collective communication.} While emerging GPU-initiated communication~\cite{nvidia_nvshmem_ibgda,deepseek_deepep_2024} optimizes small-message latency for the decoding phase, it remains uncommon in prefill deployments due to strict hardware prerequisites and marginal gains for large-volume transfers. Nevertheless, even in such environments, \sys ensures compatibility via CPU-assisted variants (e.g., host-side doorbells), while this method may introduce negligible latency overhead compared to a fully GPU-initiated execution \cite{nvidia_magnum_nvshmem3}.

\section{Related Work.}

\parab{Flow and co-flow scheduling.}
Traditional flow scheduling optimizes per-flow metrics such as average flow completion time~\cite{dctcp,dcqcn,pfabric,PDQ,pias,li2017rate,li2024flow} or deadline satisfaction~\cite{D3,d2tcp,PDQ,zhang2015guaranteeing} using policies like SJF and EDF.
Coflow scheduling~\cite{varys,Aalo,zhang2016coda,zhao2015rapier,susanto2016stream,agarwal2018sincronia,wan2025coflow} extends this abstraction by grouping related flows and optimizing Coflow Completion Time for data-parallel workloads, while mixed-flow scheduling~\cite{karuna} considers the coexistence of deadline-sensitive and best-effort traffic. In contrast, \sys{} explicitly accounts for multi-stage dependencies to schedule flows toward end-to-end TTFT attainments.

\parab{Optimization for individual stages.} Prior research focuses on optimizing specific communication phases in isolation. For collective communication, efforts improve performance by algorithms synthesis~\cite{shah2023taccl,cao2025syccl}, resolving  potential contention~\cite{cao2024crux,rajasekaran2024cassini}, or enhancing overlap through system-level pipelining~\cite{comet,zheng2025triton,hong2025flashoverlap},. In parallel, other works optimize KV cache transfers by hiding latency via prefetching~\cite{mooncakefast,splitwise}, improving efficiency via block coalescing~\cite{li2025flowkv,chen2024kvdirect} and multi-path transmission~\cite{wu2026dualpath}, or reducing data volume by trading off model quality~\cite{cachegen}. However, these approaches treat each stage in isolation, overlooking the potential contention between collective communication and KV cache transfers.

\parab{LLM serving systems.} Recent advancements in LLM serving have evolved along two complementary dimensions. Architecturally, systems adopt prefill-decode disaggregation~\cite{distserve,splitwise} and KV-cache reuse~\cite{mooncakefast,cachedattention} to improve throughput and avoid redundant computation. On the orchestration front, numerous request schedulers~\cite{thunderserve,stojkovic2025dynamollm,chen2025slos,hongsola,tempo2025} have been proposed to maximize SLO attainment. However, these systems primarily focus on optimizing compute efficiency under ideal network assumptions. \sys is orthogonal to these approaches. It bridges the gap by addressing the network contention that they overlook.

\section{Conclusion}
This paper presented \sys{}, a holistic multi-stage flow scheduling framework for LLM serving that targets end-to-end TTFT SLO attainment. \sys{} is built on the observation that request-level deadlines progressively materialize into explicit flow level deadline as prefill execution. At its core, \sys{} realizes a Defer-and-Promote scheduling principle that approximates Least-Laxity-First behavior without requiring precise laxity estimation. Through a prototype implementation and extensive evaluation, we show that \sys{} consistently improves TTFT SLO attainment across diverse models and workloads.
